

\begin{define}{Permutations without Repetition}
Arranging $r$ distinct objects chosen from $n$ distinct objects, using each at most once. 
Order matters, no object may repeat.

$$P(n,r) = \frac{n!}{(n - r)!}$$
where $0 \leq r \leq n$.
\end{define}

\textbf{What?} You have $n$ distinct items and want to choose an ordered 
list of length $r$, with $0 \leq r \leq n$, using each item at most once.

$$P(n,r) = \frac{n!}{(n - r)!}$$

\textbf{Derivation.} Think in two steps:
\begin{enumerate}
\item First choose an ordered sequence of length $n$ from all $n$ items: $n!$ ways.
\item Then "forget" the last $n-r$ positions: that overcounts by a factor of $(n - r)!$. 
Hence divide.
\end{enumerate}

\textbf{Example.} $n=5,r=3$: how many ways to pick 3-letter "words" from 
$\{A,B,C,D,E\}$ with no repeats?
$$P(5,3)=\frac{5!}{2!}=60$$

\newpage

\begin{define}{Permutations of a Multiset}
Arranging $n$ total objects when there are $n_1$ indistinguishable of type 1, 
$n_2$ indistinguishable of type 2, ..., $n_k$ of type $k$. Order matters, 
but objects of the same type are identical.

$$\frac{n!}{n_1! \cdot n_2! \cdot \cdots \cdot n_k!}$$
where $n = n_1 + n_2 + \cdots + n_k$.
\end{define}

\textbf{What?} You have $n$ total objects, but they come in $k$ types, 
with $n_1$ copies of type 1, ..., $n_k$ copies of type $k$. Objects of 
the same type are identical. You want to arrange all $n$ in a line.

$$\frac{n!}{n_1! \cdot n_2! \cdot \cdots \cdot n_k!}$$
where $n = n_1 + n_2 + \cdots + n_k$.

\textbf{Derivation.}
\begin{itemize}
\item If all $n$ were distinct, there'd be $n!$ orders.
\item But swapping among the $n_i$ copies of type $i$ doesn't change the 
arrangement, so divide by each $n_i!$.
\end{itemize}

\textbf{Example.} Arrange the letters in "BALLOON":
$$\frac{7!}{2! \cdot 2! \cdot 1! \cdot 1! \cdot 1!} = \frac{5040}{4} = 1260$$
Here $n=7$, with $n_B=1,n_A=1,n_L=2,n_O=2,n_N=1$.

\newpage

\begin{define}{Combinations without Repetition}
Selecting $r$ objects from $n$ distinct objects, where order does not 
matter and no object may repeat.

$$\binom{n}{r} = \frac{n!}{r!(n - r)!}$$
where $0 \leq r \leq n$.
\end{define}

\textbf{What?} From $n$ distinct objects, choose an unordered subset 
of size $r$ (each at most once).

$$\binom{n}{r} = \frac{n!}{r!(n - r)!}$$

\textbf{Derivation.}
\begin{itemize}
\item First count ordered sequences of length $r$: $P(n,r)=n!/(n-r)!$.
\item But order within the chosen $r$ doesn't matter, so divide by $r!$.
\end{itemize}

\textbf{Example.} From $\{1,2,3,4,5\}$, the number of 3-element 
subsets is $\binom{5}{3}=10$.

\newpage

\begin{define}{Combinations of a Multiset (with Repetition)}
Choosing $r$ objects from $n$ types, where each type may appear any 
number of times and order does not matter. Equivalently, the number 
of nonnegative integer solutions to
$$x_1 + x_2 + \cdots + x_n = r$$
where $x_i \geq 0$.

$$\binom{r + n - 1}{r} = \binom{r + n - 1}{n - 1}$$
\end{define}

\textbf{What?} You have $n$ types of object, and want to select $r$ 
items in total; each type may be chosen any number of times. Order does not matter.

$$\binom{r + n - 1}{r} = \binom{r + n - 1}{n - 1}$$

\textbf{Derivation (Stars \& Bars).}
\begin{itemize}
\item Represent your $r$ chosen items as $r$ stars in a row.
\item To group them into $n$ types, place $n-1$ bars between stars to 
mark boundaries.
\item There are $r + n - 1$ positions in total; choose which $n-1$ are 
bars (or equivalently which $r$ are stars).
\end{itemize}

\textbf{Example.} Distribute 5 identical candies among 3 kids (kids 
can get zero): $\binom{5+3-1}{5}=\binom{7}{5}=21$ ways.

\newpage

\begin{define}{Weak Compositions}
The number of ways to write a nonnegative integer $r$ as an ordered 
sum of $n$ nonnegative parts,
$$x_1 + x_2 + \cdots + x_n = r$$
where $x_i \geq 0$,
is
$$\binom{r + n - 1}{r}$$
\end{define}

\textbf{What?} Count the ordered nonnegative solutions $(x_1,\dots,x_n)$ to
$$x_1 + x_2 + \cdots + x_n = r$$
where $x_i \geq 0$.

$$\#\{\text{solutions}\} = \binom{r + n - 1}{r}$$

This is algebraically the same as "combinations of a multiset," but 
viewed as an equation in nonnegative integers.